\section{Results}
\subsection{Experiment 1}
Through experiments and observations from images, it can be seen that when the infrared sensor is placed on a white surface, the infrared light is effectively reflected after being emitted and can be received. 
At this time, the module outputs a low voltage level, and the output indicator light is illuminated. 
However, when the infrared sensor is placed on a black surface, black objects almost do not reflect infrared light. 
Therefore, when the line-following module detects a black line, it outputs a high voltage level, and the output indicator light is turned off. 
Additionally, when the line-following module is positioned too high above the ground, its detection effect is similar to that when encountering a black line. 
It will also output a high voltage level, and the indicator light will be turned off. 

We tested three sets of values at different horizontal heights of the sensor above the ground and calculated their average values as the final results~\ref{tab:values}. 
The relationship between the detection values and the height was ultimately represented by a graph~\ref{fig:graph}.

\begin{table}[H]
  \centering
  \newcolumntype{Y}{>{\centering\arraybackslash}X}
  \rowcolors{2}{gray!20}{white}
  \begin{tabularx}{0.8\textwidth}{|Y|Y|Y|Y|Y|}
    \hline
    \rowcolor{blue!30} \textbf{Altitude}/mm & & \textbf{Data} & & \textbf{Average} \\ 
    \hline
     5 & 33 & 34 & 33 & 33.3 \\ 
    \hline
    10 & 37 & 38 & 38 & 37.7 \\ 
    \hline
    15 & 42 & 43 & 42 & 42.3 \\ 
    \hline
    20 & 48 & 49 & 48 & 48.3 \\ 
    \hline
    25 & 59 & 54 & 58 & 57.0 \\ 
    \hline
    30 & 65 & 70 & 86 & 73.7 \\ 
    \hline
    35 & 294 & 275 & 315 & 294.7 \\ 
    \hline
    40 & 560 & 465 & 510 & 511.7 \\ 
    \hline
    45 & 765 & 680 & 648 & 697.7 \\ 
    \hline
    50 & 895 & 800 & 846 & 847.0 \\ 
    \hline
  \end{tabularx}
  \caption{Data and Average Values}
  \label{tab:values}
\end{table}

\begin{figure}[H]
  \centering
	\includegraphics[width=0.8\linewidth]{../../Photos/lab4_graph}
  \caption{Altitude - Sensor Value Graph}
  \label{fig:graph}
\end{figure}

Based on the data above, it can be seen that when the sensor is 3 cm above the ground, the sensor values undergo a sudden change. 
Therefore, in order to ensure stability, the sensor must be positioned no more than 3 cm above the ground.

\subsection{Experiment 2}
Based on the aforementioned experiment, we first recorded the sensor readings. 
Subsequently, we established a recognition criterion: 
when the sensor reading is less than 100, it is marked as “0”; 
and when the reading approaches 1023, it is marked as “1”. 
Here, “0” indicates that the sensor has detected a white surface, while “1” indicates that it has detected a black line. 
Under these conditions, the sensor readings under different scenarios are shown in the figures below ():



% Add all appropriate figures and analysis.
% No formal word count for this section, however
% please keep text as concise as possible.